\section{Experimental design}\label{sec:exp}
\textbf{Graphs.} Synthetic ($n{=}3000$ unless stated; largest component): Erd\H{o}s--R\'enyi (ER, mean degree 16), Barab\'asi--Albert (BA, $m{=}8$), Watts--Strogatz (WS, $k{=}16,\beta{=}0.1$), stochastic block model (SBM, 30 blocks, mean degree 16, fraction of between-block edges $\texttt{mix}$ default $0.1$); public graphs: SNAP ego-Facebook \citep{mcauley2012} (4039 nodes, mean degree 43.7) and Facebook100 Stanford3 \citep{traud2012,rossi2015} (11{,}586 nodes, mean degree 98.1).

\textbf{Outcome model.} $\mu_i(z)=a_i+\tau_dz_i+\gamma\,g_i(z)$ with $\tau_d{=}\gamma{=}1$ and $\sigma{=}1$ by default. Baselines $a_i=1+0.5\,\mathrm{zlogdeg}_i+0.5\,b_{c(i)}+0.7\,\nu_i$ ($b_c$: Louvain-community effect, $\nu_i$ i.i.d.\ $N(0,1)$) make outcomes homophilous and degree-dependent and are \emph{not} fully known to the regression. Writing $x_i=\frac1{d_i}\sum_jA_{ij}z_j$, the eight mechanisms differ in $g$, normalised so that $g(\mathbf 1)-g(\mathbf 0)=1$ (GATE$=2$, so relative errors are comparable): \texttt{linfrac}: $x_i$; \texttt{lincount}: $\sum_jA_{ij}z_j/\bar d$ (hub amplification); \texttt{sat}: $(1-e^{-4x_i})/(1-e^{-4})$ (diminishing returns); \texttt{quad}: $x_i^2$; \texttt{thresh}: logistic in $x_i$ (centre $0.7$, slope $10$, rescaled; complex contagion); \texttt{twohop}: $\frac12x_i+\frac12(Wx)_i$; \texttt{hetero}: $(1+0.5u_i)x_i$ with latent $u_i$ partly correlated with degree; \texttt{peer}: equilibrium peer effects $\mu=(I-\rho W)^{-1}(a+\tau_dz)$, $\rho{=}0.5$, $W$ row-normalised (interference of unbounded range). The GATE of every population is computed exactly from $\mu(\mathbf 1)-\mu(\mathbf 0)$.

\textbf{Designs.} \emph{Bernoulli}($p$), default $p{=}0.5$; \emph{cluster}: Louvain clusters \citep{blondel2008} ($K$ clusters, fraction cut edges reported), half of the clusters treated by complete randomisation; \emph{randomised saturation} (clusters draw $p_c\in\{.1,.3,.5,.7,.9\}$, units Bernoulli($p_c$)).

\textbf{Estimators.} \textbf{DM}; OLS on $(z,x)$ (\textbf{OLS-x}, the model of Prop.~\ref{prop:bern}), plus degree (\textbf{xd}), plus $x^2$ and $zx$ (\textbf{quad}), plus two-hop exposure (\textbf{x2}), a 10-feature polynomial/interaction model (\textbf{rich}), count exposure without/with degree control (\textbf{c}, \textbf{cd}); \textbf{GBDT}: histogram gradient boosting on $(z,x,\text{count},\text{2-hop},\text{degree})$ predicting at all-1/all-0 features; \textbf{H\'ajek}: exposure-conditional H\'ajek estimator for ``treated with $x_i\ge0.75$'' vs.\ ``control with $x_i\le0.25$'' with exposure probabilities obtained by Monte-Carlo ($M{=}400$ design draws); and an \textbf{oracle} OLS using the true exposure basis (plus degree). H\'ajek targets a coarsened ($q{=}0.75$) contrast and so is biased for the GATE by construction.

\textbf{Metrics and uncertainty.} Per configuration we run $R{=}100$--$500$ independent assignments and noise draws (fixed population) and report \emph{bias}, \emph{SD}, and \emph{RMSE} (primary: it is the quantity a practitioner experiences, and it penalises bias and variance), mostly as relative RMSE $=\mathrm{RMSE}/|\tau|$; 95\% bootstrap CIs over replicates for RMSE and Monte-Carlo standard errors for bias. For each replicate we also refit on noiseless $\mu(z)$ to separate $S$ (assignment-induced; its mean is the \emph{bias}) from $N$ (noise) as in Prop.~\ref{prop:decomp}. Coverage of 95\% intervals is evaluated in \S\ref{sec:res}.
